\begin{table}[htbp]
  \centering
  \caption{Administrative states and distance to metropolitan area (in Km)}
  \begin{tabular}{lrr}
    \midrule
    \midrule
    \multicolumn{1}{c}{\textbf{State}} &
    \multicolumn{1}{c}{\textbf{Distance to metro area}} &
    \multicolumn{1}{c}{\textbf{Potential pollution source?}} \\
    \multicolumn{1}{c}{} &
    \multicolumn{1}{c}{\textbf{(km)}} &
    \multicolumn{1}{c}{\textbf{($\leq 20$ km)}} \\
    \midrule
    São Paulo & 0.00 & 1 \\
    Minas Gerais & 33.41 & 0 \\
    Rio de Janeiro & 95.22 & 0 \\
    Paraná & 153.35 & 0 \\
    Santa Catarina & 260.19 & 0 \\
    Espírito Santo & 494.11 & 0 \\
    Mato Grosso do Sul & 541.21 & 0 \\
    Goiás & 544.55 & 0 \\
    Rio Grande do Sul & 557.82 & 0 \\
    Distrito Federal & 798.02 & 0 \\
    Bahia & 841.44 & 0 \\
    Mato Grosso & 863.68 & 0 \\
    Tocantins & 1,097.77 & 0 \\
    Piauí & 1,365.56 & 0 \\
    Maranhão & 1,423.98 & 0 \\
    Pará & 1,531.08 & 0 \\
    Sergipe & 1,573.01 & 0 \\
    Pernambuco & 1,628.52 & 0 \\
    Alagoas & 1,750.51 & 0 \\
    Rondônia & 1,808.51 & 0 \\
    Ceará & 1,852.00 & 0 \\
    Paraíba & 1,896.80 & 0 \\
    Rio Grande do Norte & 2,032.35 & 0 \\
    Amazonas & 2,049.49 & 0 \\
    Amapá & 2,510.94 & 0 \\
    Acre & 2,626.44 & 0 \\
    Roraima & 2,926.29 & 0 \\
    \bottomrule
    \bottomrule
  \end{tabular}
  \label{table_state_metro_distances}
\end{table}
